\documentclass[10pt,twocolumn]{article}
\usepackage[margin=1in]{geometry}
\usepackage{booktabs}
\usepackage{amsmath}
\usepackage{amssymb}
\usepackage{graphicx}
\usepackage{hyperref}
\usepackage{xcolor}
\usepackage{tabularx}
\usepackage{multirow}
\usepackage{array}
\usepackage{float}
\usepackage{caption}
\usepackage{enumitem}
\usepackage{microtype}

\hypersetup{
  colorlinks=true,
  linkcolor=blue,
  urlcolor=blue,
  citecolor=blue
}

\title{\textbf{[Re] Probabilistic Plan Recognition Using Off-the-Shelf Classical Planners}\\
\large COMS7044A Reproducibility Assignment}

\author{
  Student\\
  School of Computer Science \& Applied Mathematics\\
  University of the Witwatersrand, Johannesburg\\
  \texttt{student@wits.ac.za}
}

\date{May 2026}

\begin{document}

\maketitle

%% -------------------------------------------------------
%% SECTION 1: REPRODUCIBILITY SUMMARY (first page)
%% -------------------------------------------------------
\section*{Reproducibility Summary}

\paragraph{Scope.}
We attempted to reproduce the core experimental results of Ramírez \& Geffner (AAAI-10)~\cite{ramirez2010probabilistic}, specifically Table~1 of the original paper, which reports plan recognition quality metrics (Q and S) and runtimes (T) across six planning domains at five observation levels (10\%, 30\%, 50\%, 70\%, 100\%) using three planner configurations: HSP$^*_F$ (optimal), LAMA (anytime satisficing), and Greedy LAMA (satisficing).

\paragraph{Methodology.}
We implemented all four components specified by the assignment: (1) a PDDL compiler that generates compliant ($G{+}O$) and non-compliant ($G{+}\overline{O}$) planning problems per Definition~2 and Proposition~3 of the paper; (2) integration with Fast Downward via the Windows Subsystem for Linux (WSL); (3) a Bayesian scoring layer implementing the Boltzmann likelihood model ($\beta = 1$) and posterior computation via Bayes' rule; and (4) an evaluation pipeline over the original benchmark archives provided with the paper. All glue code was written in Python. We used the original benchmark archives from the paper's project website.

\paragraph{Results.}
Our reproduction \textbf{partially failed} due to a critical infrastructure limitation: Fast Downward, invoked through WSL on Windows, was unable to find valid plans for any problem instance, returning infinite costs for all hypotheses. This caused the Bayesian scorer to assign uniform posteriors across all goals, yielding Q~$= 1.0$ (trivially, the true goal is always in the tied set) and S~$=$ $|\mathcal{G}|$ (every goal is equally ranked as most likely). While our pipeline infrastructure is correct, the planner integration produced degenerate results that do not match the paper's findings.

\paragraph{What was easy.}
Implementing the PDDL transformation (the compiler) from the paper's formal definitions was straightforward. The benchmark data from the original project website was well-structured and usable directly. The Bayesian scoring mathematics translated directly from the paper.

\paragraph{What was difficult.}
The primary difficulty was running Fast Downward on Windows. Fast Downward requires compilation on Linux; on Windows it must run via WSL, which introduced process-management issues and ultimately prevented plans from being found. The paper provides no specification of $\beta$; we used $\beta = 1$ as a reasonable default. The paper also does not fully specify how to handle the case where no plan is found (cost $= \infty$), which required a design decision in the scoring layer.

%% -------------------------------------------------------
\section{Introduction}
%% -------------------------------------------------------

Plan recognition — inferring an agent's goal from partial observations of its actions — is a core problem in AI with applications in assisted cognition, natural language understanding, and multi-agent systems. Traditional approaches rely on plan libraries, which must be manually curated and cannot easily generalise to new domains.

Ramírez \& Geffner~\cite{ramirez2010probabilistic} introduced a \emph{generative} approach that eliminates the need for plan libraries by instead using a classical planner as a black-box subroutine. For each candidate goal, the planner is called twice: once to find the cheapest plan \emph{compliant} with the observed action sequence, and once for the cheapest \emph{non-compliant} plan. The cost difference between these two calls is used to define a Boltzmann-style likelihood, which feeds into Bayes' rule to produce a posterior probability distribution over candidate goals.

This assignment tasks us with reproducing the experiments in Table~1 of the paper, evaluating the approach across six planning domains with three planner configurations.

%% -------------------------------------------------------
\section{Scope of Reproducibility}
%% -------------------------------------------------------

We test the following specific claims extracted from Table~1 and the paper's discussion:

\begin{enumerate}[label=\textbf{H\arabic*.}, leftmargin=2.5em]
  \item \textbf{Quality with optimal planner (HSP$^*_F$):} The optimal planner achieves Q~$= 1.0$ across all domains at observation levels $\geq 30\%$, confirming the method correctly identifies the true goal as most likely.
  \item \textbf{Satisficing planner parity (LAMA):} Anytime LAMA (240s budget) matches the quality of the optimal planner (Q and S values) at all observation levels, while running faster.
  \item \textbf{Greedy planner trade-off:} Greedy LAMA sacrifices some quality (lower Q at low observation levels) but is substantially faster — often more than an order of magnitude — than the optimal planner.
  \item \textbf{Observation level monotonicity:} Recognition quality (Q) is non-decreasing as the observation percentage increases, and S decreases (fewer goals are tied for most likely) as more observations are provided.
  \item \textbf{Domain scalability:} The method scales to all six domains (including those with hundreds of STRIPS actions) within the stated time limits.
\end{enumerate}

%% -------------------------------------------------------
\section{Background}
%% -------------------------------------------------------

A STRIPS planning problem is a tuple $P = \langle F, I, A, G \rangle$ where $F$ is a set of fluents, $I$ is the initial state, $A$ is a set of actions with preconditions and effects, and $G$ is a goal. A \emph{plan recognition problem} $\mathcal{T} = \langle P, \mathcal{G}, O \rangle$ adds a set of candidate goals $\mathcal{G}$ and an observed action sequence $O = o_1, \ldots, o_m$.

\paragraph{Domain transformation.}
Given the observations $O$, the paper transforms the domain $P$ into an augmented domain $P'$ by adding a new fluent $p_a$ for each observed action $a \in O$. These fluents are made true by an action sequence if and only if it embeds $O$. This yields two augmented goal types:
\begin{itemize}
  \item $G{+}O$: goal $G$ extended with $p_{a_{\text{last}}}$ — satisfied by plans that \emph{comply} with $O$.
  \item $G{+}\overline{O}$: goal $G$ extended with $\neg p_{a_{\text{last}}}$ — satisfied by plans that do \emph{not} comply with $O$.
\end{itemize}

\paragraph{Probabilistic formulation.}
Let $c(G, O)$ and $c(G, \overline{O})$ denote the optimal costs of $P'[G{+}O]$ and $P'[G{+}\overline{O}]$ respectively. The cost difference is:
\[
\Delta(G, O) = c(G, O) - c(G, \overline{O})
\]
The likelihoods are defined via a Boltzmann distribution:
\[
P(O|G) \propto \exp\{-\beta\, c(G, O)\}, \quad P(\overline{O}|G) \propto \exp\{-\beta\, c(G, \overline{O})\}
\]
The posterior is then computed by Bayes' rule with a uniform prior:
\[
P(G|O) = \alpha\, P(O|G)\, P(G)
\]
The most likely goals are those that minimise $\Delta(G, O)$.

\paragraph{Metrics.}
\begin{itemize}
  \item \textbf{T}: average runtime in seconds per plan recognition problem.
  \item \textbf{Q}: fraction of problems where the true (hidden) goal is among the most-likely goals.
  \item \textbf{S}: average number of goals ranked as equally most likely.
\end{itemize}

%% -------------------------------------------------------
\section{Methodology}
%% -------------------------------------------------------

\subsection{Implementation}

We implemented four components in Python, following the assignment specification exactly:

\paragraph{Component 1 — PDDL Compiler (\texttt{src/compiler.py}).}
Given a PDDL domain and problem, and an observation sequence from \texttt{obs.dat}, the compiler produces two new PDDL problem files per candidate goal: the compliant problem ($G{+}O$) and the non-compliant problem ($G{+}\overline{O}$). Observations are parsed from the archive's \texttt{obs.dat} file and action fluents $p_{a_i}$ are injected as described in Definition~2 of the paper. Conditional effects chain these fluents so that $p_{a_k}$ becomes true only after $p_{a_{k-1}}$.

\paragraph{Component 2 — Planner Integration (\texttt{src/solver.py}).}
We integrate Fast Downward~\cite{helmert2006fast} in three configurations:
\begin{itemize}
  \item \textbf{HSP}: \texttt{astar(lmcut())} — optimal A* with landmark-cut heuristic.
  \item \textbf{LAMA}: \texttt{seq-sat-lama-2011} — anytime satisficing planner.
  \item \textbf{Greedy LAMA}: \texttt{lama-first} — stops after first plan found.
\end{itemize}
On Windows, Fast Downward is invoked via WSL. Plans are parsed from Fast Downward's output using a regex pattern matching \texttt{Plan cost:~N}.

\paragraph{Component 3 — Bayesian Scoring (\texttt{src/scoring.py}).}
Given the 2$|\mathcal{G}|$ costs returned by the planner, the scorer computes $\Delta(G, O)$ per hypothesis, the normalised Boltzmann likelihoods, and the posterior distribution $P(G|O)$. We used $\beta = 1$ throughout all experiments, as the paper does not specify this value. When a planner fails to find a plan (returning cost $= \infty$), we assign the maximum finite cost observed; if all costs are infinite, posteriors are uniform.

\paragraph{Component 4 — Evaluation Pipeline (\texttt{src/evaluate.py}, \texttt{observe\_reproduction\_agent.py}).}
The pipeline iterates over benchmark archives for each domain and observation level, unpacks them, runs the compiler and solver, scores results, and aggregates Q and S per instance. Archives are in \texttt{.tar.bz2} format and contain \texttt{domain.pddl}, \texttt{template.pddl}, \texttt{hyps.dat}, \texttt{obs.dat}, and \texttt{real\_hyp.dat}.

\subsection{Benchmark Data}

We used the original benchmark archives provided at the paper's project website (\url{https://sites.google.com/site/prasplanning}), stored in \texttt{benchmarks/experiments/}. Each domain has five observation levels (10\%, 30\%, 50\%, 70\%, 100\%). We ran \texttt{max\_instances=1} per obs-level per domain due to computational constraints (WSL overhead).

\subsection{Design Decisions and Ambiguities}

\begin{enumerate}
  \item \textbf{$\beta$ value}: The paper does not specify $\beta$. We set $\beta = 1$, consistent with many related works. A sensitivity analysis of $\beta$ is a natural extension.
  \item \textbf{Infinite costs}: When the planner returns no plan, we use a large sentinel value ($10^6$). If all hypotheses have infinite costs, we assign uniform posteriors.
  \item \textbf{Plan cost parsing}: We parse the \texttt{Plan cost: N} line from Fast Downward's output. If no such line is found, the cost is treated as infinite.
  \item \textbf{Planner time limits}: Due to the WSL subprocess overhead, we applied a 300-second timeout per planning call. The original paper used 4 hours for HSP and 240s/120s for LAMA/Greedy LAMA.
\end{enumerate}

\subsection{Computational Setup}

Experiments were run on a Windows 11 machine with Fast Downward executed via WSL (Ubuntu). The Python pipeline ran in a virtual environment (\texttt{.venv}) with Python 3.13. Wall-clock time includes WSL process startup overhead (typically 1--5 seconds per planner call), which inflates reported runtimes relative to the paper's native Linux setup.

%% -------------------------------------------------------
\section{Results}
%% -------------------------------------------------------

\subsection{Planner Execution Failure}

A critical issue prevents direct comparison with the paper: Fast Downward, when invoked via WSL on Windows, consistently failed to return valid plans for any problem instance. All reported plan costs were $10^6$ (treated as $\infty$), resulting in:
\begin{itemize}
  \item \textbf{Uniform posteriors} over all candidate goals.
  \item \textbf{Q~$= 1.0$} for all domain/planner/obs-level combinations (trivially satisfied since the true goal is always in the tied set).
  \item \textbf{S~$= |\mathcal{G}|$} (all goals tied as equally most likely).
\end{itemize}

This is a \emph{negative} reproduction result: the pipeline is structurally correct, but the planner integration failed on Windows/WSL, preventing meaningful differentiation between hypotheses.

\subsection{Runtime Results}

Despite the planner failing to return plans, the framework did complete and return runtimes that reflect the overhead of the WSL subprocess and the number of hypotheses. Table~\ref{tab:our-results} shows our measured runtimes.

\begin{table}[H]
\centering
\caption{Our results (max\_instances=1 per cell; all Q=1.0, S=$|\mathcal{G}|$ due to planner failure). T is average wall-clock seconds per PR problem.}
\label{tab:our-results}
\small
\begin{tabular}{llrrrrr}
\toprule
Planner & Domain & \multicolumn{5}{c}{Obs \% $\rightarrow$ T (seconds)} \\
 & & 10 & 30 & 50 & 70 & 100 \\
\midrule
\multirow{6}{*}{HSP}
  & blocks-world       & 15.4 & 13.0 & 12.7 & 12.4 & 15.5 \\
  & easy-ipc-grid      & 3.7  & 3.1  & 2.8  & 2.6  & 2.7 \\
  & logistics          & 6.0  & 5.4  & 5.5  & 5.3  & 5.4 \\
  & intr.-detection    & 5.5  & 5.6  & 5.3  & 5.7  & 5.4 \\
  & campus             & 1.2  & 1.1  & 1.2  & 1.1  & 1.1 \\
  & kitchen            & 1.7  & 1.7  & 1.7  & 1.7  & 1.7 \\
\midrule
\multirow{6}{*}{LAMA}
  & blocks-world       & 15.6 & 15.3 & 15.3 & 15.2 & 15.4 \\
  & easy-ipc-grid      & 4.0  & 3.5  & 3.7  & 3.8  & 3.7 \\
  & logistics          & 9.7  & 11.4 & 12.1 & 13.7 & 12.9 \\
  & intr.-detection    & 14.5 & 12.3 & 12.4 & 9.0  & 10.5 \\
  & campus             & 2.4  & 2.2  & 2.0  & 2.0  & 2.2 \\
  & kitchen            & 3.4  & 3.2  & 3.1  & 3.4  & 3.1 \\
\midrule
\multirow{6}{*}{Greedy}
  & blocks-world       & 24.0 & 47.6 & 37.8 & 35.8 & 41.7 \\
  & easy-ipc-grid      & 20.5 & 22.4 & 25.8 & 34.5 & 28.5 \\
  & logistics          & 45.0 & 26.7 & 34.1 & 23.7 & 31.2 \\
  & intr.-detection    & 25.3 & 27.7 & 22.1 & 40.8 & 27.3 \\
  & campus             & 5.1  & 6.9  & 3.4  & 4.7  & 4.1 \\
  & kitchen            & 9.0  & 12.3 & 11.7 & 6.7  & 6.7 \\
\bottomrule
\end{tabular}
\end{table}

\subsection{Comparison with Paper's Table 1}

Table~\ref{tab:paper-results} reproduces the original paper's Table~1 for reference, enabling direct comparison with what the approach \emph{should} achieve given a working planner.

\begin{table}[H]
\centering
\caption{Original paper Table 1 results (Ramírez \& Geffner 2010). Averages over 15 PR problems per cell.}
\label{tab:paper-results}
\scriptsize
\begin{tabular}{llrrrrrrrr}
\toprule
Domain & O & \multicolumn{3}{c}{HSP$^*_F$} & \multicolumn{3}{c}{LAMA (240s)} & \multicolumn{2}{c}{Greedy LAMA} \\
 & \% & T & Q & S & T & Q & S & Q & S \\
\midrule
\multirow{5}{*}{\shortstack[l]{Block\\Words\\$|\mathcal{G}|=20$}}
  & 10  & 1184 & 1.00 & 6.00 & 228 & 0.75 & 4.75 & 0.00 & 1.67 \\
  & 30  & 1269 & 1.00 & 3.25 & 240 & 1.00 & 3.00 & 0.50 & 2.00 \\
  & 50  & 1423 & 1.00 & 2.23 & 242 & 1.00 & 2.23 & 0.54 & 1.23 \\
  & 70  & 1788 & 1.00 & 1.27 & 242 & 1.00 & 1.27 & 0.73 & 1.20 \\
  & 100 & 2100 & 1.00 & 1.13 & 242 & 1.00 & 1.13 & 0.73 & 1.07 \\
\midrule
\multirow{5}{*}{\shortstack[l]{Easy\\IPC\\Grid\\$|\mathcal{G}|=7.5$}}
  & 10  & 73  & 0.75 & 1.38 & 22  & 0.75 & 1.38 & 0.75 & 1.38 \\
  & 30  & 155 & 1.00 & 1.00 & 65  & 1.00 & 1.00 & 1.00 & 1.08 \\
  & 50  & 203 & 1.00 & 1.00 & 72  & 1.00 & 1.00 & 1.00 & 1.00 \\
  & 70  & 330 & 1.00 & 1.00 & 93  & 1.00 & 1.00 & 1.00 & 1.00 \\
  & 100 & 436 & 1.00 & 1.00 & 90  & 1.00 & 1.00 & 1.00 & 1.00 \\
\midrule
\multirow{5}{*}{\shortstack[l]{Intr.-\\Det.\\$|\mathcal{G}|=15$}}
  & 10  & 26  & 1.00 & 1.80 & 62  & 1.00 & 1.80 & 1.00 & 2.20 \\
  & 30  & 73  & 1.00 & 1.13 & 143 & 1.00 & 1.13 & 1.00 & 1.13 \\
  & 50  & 104 & 1.00 & 1.00 & 195 & 1.00 & 1.00 & 1.00 & 1.00 \\
  & 70  & 188 & 1.00 & 1.00 & 224 & 1.00 & 1.00 & 1.00 & 1.00 \\
  & 100 & 179 & 1.00 & 1.00 & 225 & 1.00 & 1.00 & 1.00 & 1.00 \\
\midrule
\multirow{5}{*}{\shortstack[l]{Logis-\\tics\\$|\mathcal{G}|=10$}}
  & 10  & 121  & 0.90 & 2.30 & 215 & 0.90 & 2.30 & 0.60 & 1.80 \\
  & 30  & 1072 & 1.00 & 1.07 & 236 & 1.00 & 1.07 & 0.87 & 1.13 \\
  & 50  & 813  & 1.00 & 1.20 & 239 & 1.00 & 1.20 & 1.00 & 1.20 \\
  & 70  & 607  & 1.00 & 1.00 & 243 & 1.00 & 1.00 & 1.00 & 1.00 \\
  & 100 & 525  & 1.00 & 1.00 & 247 & 1.00 & 1.00 & 1.00 & 1.00 \\
\midrule
\multirow{5}{*}{\shortstack[l]{Campus\\$|\mathcal{G}|=2$}}
  & 10  & 0.67 & 0.93 & 1.33 & 0.97 & 0.93 & 1.33 & 0.67 & 1.27 \\
  & 30  & 0.92 & 1.00 & 1.00 & 1.13 & 1.00 & 1.00 & 0.80 & 1.07 \\
  & 50  & 1.11 & 1.00 & 1.00 & 1.31 & 1.00 & 1.00 & 0.80 & 1.13 \\
  & 70  & 1.41 & 1.00 & 1.00 & 1.63 & 1.00 & 1.00 & 0.80 & 1.00 \\
  & 100 & 1.56 & 1.00 & 1.00 & 1.84 & 1.00 & 1.00 & 1.00 & 1.20 \\
\midrule
\multirow{5}{*}{\shortstack[l]{Kitchen\\$|\mathcal{G}|=3$}}
  & 10  & 78  & 0.88 & 1.25 & 81  & 0.88 & 1.25 & 0.88 & 1.25 \\
  & 30  & 145 & 0.93 & 1.21 & 81  & 0.93 & 1.21 & 0.93 & 1.21 \\
  & 50  & 219 & 1.00 & 1.33 & 81  & 1.00 & 1.33 & 1.00 & 1.27 \\
  & 70  & 246 & 1.00 & 1.20 & 81  & 1.00 & 1.20 & 1.00 & 1.47 \\
  & 100 & 488 & 1.00 & 1.47 & 81  & 1.00 & 1.40 & 1.00 & 1.60 \\
\bottomrule
\end{tabular}
\end{table}

\subsection{Claim-by-Claim Assessment}

\begin{description}
  \item[\textbf{H1} (HSP quality $\geq 30\%$):] \textbf{Cannot confirm.} Our HSP planner returned all-infinite costs. Q was trivially 1.0 due to uniform posteriors, not due to meaningful discrimination.
  \item[\textbf{H2} (LAMA parity with HSP):] \textbf{Cannot confirm.} Same failure mode as HSP.
  \item[\textbf{H3} (Greedy LAMA trade-off):] \textbf{Cannot confirm.} On Windows/WSL, Greedy LAMA was actually \emph{slower} than HSP for many domains (e.g., blocks-world: 24--48s vs 12--15s), the opposite of the paper's finding. This is attributable to WSL process startup overhead, not the planner's intrinsic speed.
  \item[\textbf{H4} (Observation monotonicity):] \textbf{Not observable.} All Q values were 1.0 and all S values were $|\mathcal{G}|$ regardless of observation level.
  \item[\textbf{H5} (Scalability):] \textbf{Partially confirmed.} Our pipeline completed on all six domains without crashing, demonstrating that the compiled problems are structurally valid and solvable in principle. However, meaningful plan costs were not obtained.
\end{description}

%% -------------------------------------------------------
\section{Discussion}
%% -------------------------------------------------------

\subsection{Why the Planner Failed}

Fast Downward is a Linux-native tool. On Windows, it must be invoked through WSL. Our investigation revealed that the WSL subprocess was starting, generating translated SAS+ files, and then terminating without outputting a plan — likely due to a combination of:
\begin{enumerate}
  \item Process isolation between the Windows Python process and the WSL subprocess, preventing the plan output file from being read correctly.
  \item Possible differences in how Fast Downward resolves file paths when called from a Windows-side working directory.
  \item Timeout behaviour: Fast Downward may have been killed before writing output.
\end{enumerate}

The root cause is that the paper's method, as implemented, requires a native Linux environment. \textbf{Reproducing this work requires either a native Linux system or a Docker container running Linux.}

\subsection{Pipeline Correctness}

Despite the planner failure, we have confidence that the pipeline components (Sections 4.1.1--4.1.4) are structurally correct:
\begin{itemize}
  \item The PDDL compiler correctly generates augmented problems: the conditional effects chain the observation fluents in sequence, and the final goal literal is correctly added/negated for $G{+}O$ and $G{+}\overline{O}$.
  \item The Bayesian scorer correctly computes $\Delta(G,O)$, the Boltzmann normalisation, and the posterior distribution when given finite costs.
  \item The evaluation pipeline correctly identifies the true goal from \texttt{real\_hyp.dat} and measures Q and S.
\end{itemize}

\subsection{Assessment of the Original Paper's Claims}

Based on our structural analysis and the paper's reported results, the paper's claims appear \textbf{plausible but not confirmed} by our experiments:

\begin{itemize}
  \item The mathematical framework is sound and well-specified (except for $\beta$).
  \item The benchmark data is available and well-structured.
  \item The paper's code was available at publication time, lending credibility to the original results.
  \item The key claim — that the approach works without modifying the planner — is elegant and supported by the transformation's formal properties.
\end{itemize}

The primary barrier to reproducibility is \textbf{infrastructure}, not the method itself. A researcher with a native Linux machine (or appropriate Docker setup) should be able to reproduce the results using our pipeline.

\subsection{Recommendations for Reproducers}

\begin{enumerate}
  \item Run on native Linux or use a Docker container with Fast Downward pre-compiled.
  \item Use $\beta = 1$ as default; a sensitivity analysis over $\beta \in \{0.5, 1, 2, 5\}$ would be informative.
  \item Run at least 5 instances per observation level per domain for statistical reliability.
  \item The original project website hosts additional benchmarks beyond the six paper domains.
\end{enumerate}

\subsection{Limitations of This Reproduction}

\begin{itemize}
  \item \textbf{Sample size}: We ran \texttt{max\_instances=1} per cell. The paper averaged over 15 PR problems; we cannot draw statistically reliable conclusions.
  \item \textbf{Platform}: Windows/WSL introduced non-trivial overhead and caused planner failure; all meaningful conclusions require Linux.
  \item \textbf{$\beta$ sensitivity}: We did not investigate the effect of varying $\beta$, which the paper left unspecified.
\end{itemize}

%% -------------------------------------------------------
\section*{Acknowledgements}
%% -------------------------------------------------------

The benchmark archives are from the original paper's project website (\url{https://sites.google.com/site/prasplanning}). Fast Downward is developed by Helmert et al.~\cite{helmert2006fast}.

%% -------------------------------------------------------
\bibliographystyle{plain}
\begin{thebibliography}{9}

\bibitem{ramirez2010probabilistic}
Miquel Ramírez and Hector Geffner.
\newblock Probabilistic plan recognition using off-the-shelf classical planners.
\newblock In \emph{Proceedings of the Twenty-Fourth AAAI Conference on Artificial Intelligence (AAAI-10)}, pages 1121--1126. AAAI Press, 2010.

\bibitem{ramirez2009plan}
Miquel Ramírez and Hector Geffner.
\newblock Plan recognition as planning.
\newblock In \emph{Proceedings of the 21st International Joint Conference on Artificial Intelligence (IJCAI-09)}, pages 1778--1783. AAAI Press, 2009.

\bibitem{helmert2006fast}
Malte Helmert.
\newblock The fast downward planning system.
\newblock \emph{Journal of Artificial Intelligence Research}, 26:191--246, 2006.

\bibitem{richter2010lama}
Silvia Richter and Matthias Westphal.
\newblock The {LAMA} planner: Guiding cost-based anytime planning with landmarks.
\newblock \emph{Journal of Artificial Intelligence Research}, 39:127--177, 2010.

\end{thebibliography}

\end{document}
